\documentclass{report}
\usepackage{amsmath,amssymb}
\setlength{\parindent}{0mm}
\setlength{\parskip}{1em}
\begin{document}
\begin{center}
\rule{6in}{1pt} \
{\large Arvind Krishna Saibaba \\
{\bf Fast Iterative Methods for Bayesian Inverse Problems}}

2311 Stinson Drive 3126 \\ Raleigh \\ NC 27606
\\
{\tt asaibab@ncsu.edu}\\
Julianne Chung\end{center}

Inverse problems arise in various scientific applications, and Bayesian
approaches are often used for solving statistical inverse problems, where
unknowns are modeled as random fields and Bayes' rule {is used} to infer
unknown parameters by conditioning on measurements. In this work, we
develop iterative methods for solving the following least squares
problem, which is equivalent to computing the maximum a posteriori (MAP)
estimator
\begin{equation}\label{eq_map}
\hat{x}_\text{MAP} \equiv {\arg}\min_x - \log p(x|d) \> =\>{\arg}\min_x
\> \frac{1}{2}\|d-Ax\|_{R^{-1}}^2 +
\frac{\lambda^2}{2}\|x-\mu\|_{Q^{-1}}^2.
\end{equation}

We model the prior covariance matrix $Q$ to have entries of the form
$\kappa(\vec{x}_i,\vec{x}_j)$, where $\kappa(\cdot,\cdot)$ is a
covariance kernel, and $\{\vec{x}_i\}_{i=1}^n$ are spatial points
representing the discrete image. We choose $\kappa$ from the Mat\'{e}rn
class of covariance kernels, which not only represent a rich class of
priors but also
are convenient from a theoretical and modeling point of view. However,
one of the main challenges of using the Mat\'{e}rn kernels is that
$Q^{-1}$ is difficult to obtain and work with, whereas matrix-vector
products (matvecs) with $Q$ can be handled efficiently in
$\mathcal{O}(n\log n)$ using FFT methods (on regular grids) and using
$\mathcal{H}$-matrices (on irregular grids). In our approach, we make an
appropriate change of variables and consider a hybrid Golub-Kahan
bidiagonalization iterative process with weighted inner products $\langle
\cdot, \cdot \rangle_{R^{-1}} $ and $\langle \cdot, \cdot \rangle_{Q} $.
Our approach avoids forming (either explicitly, or matvecs with) the
square root or inverse of the prior covariance matrix $Q$. The resulting
algorithms are efficient and scalable to large problem sizes.

Our proposed hybrid approach can automatically determine regularization
parameter $\lambda$, which controls the relative importance between the
prior and the data-misfit part of the likelihood. Hybrid methods project
the original least squares problem onto a smaller dimensional problem,
using orthonormal bases for the Krylov subspaces generated during the
iterative process, and regularization parameters are sought by optimizing
an appropriate functional on the projected problem. In this talk, we show
how several commonly used functionals can be appropriately modified to
handle weighted norms, as in Equation~\eqref{eq_map}. Since the
regularization parameters are determined on-the-fly during the iterative
process, this approach is attractive for large-scale inverse problems.


In order to characterize the uncertainty in parameter reconstruction, we
must fully specify the posterior distribution characterized by the MAP
estimator and the posterior covariance matrix. However, the posterior
covariance matrix is dense, and therefore, storing and computing it is
infeasible for large-scale problems. Using the bidiagonalization process,
we describe an efficient approach to generate approximations to useful
uncertainty measures based on the posterior covariance matrix. In
particular, we estimate the variance, which are the diagonals of the
posterior covariance, and show how to generate conditional realizations,
which are samples from the posterior distribution.

We will demonstrate the benefits of our proposed algorithms on
challenging 1D and 2D image reconstruction problems.






\end{document}
